Entre cargas, para cantidades en mol y tiempo en min,
\[
a'=-ab,\qquad b'=-ab,\qquad p'=ab.
\]
En cada carga sólo \(b\) salta por la dosis añadida. Los resultados son
\[
\begin{array}{c|cc|c}
& p(1)\ (\mathrm{mol})&\max_{[0,1]}b\ (\mathrm{mol})&\text{cumple ambos criterios}\\\hline
\mathrm I&1.675052\ldots&2&\text{no}\\
\mathrm{II}&1.483806\ldots&1.279530\ldots&\text{sí}
\end{array}
\]
Con \(q=e^{-1}\), sea \(b_*=1+2q/(3-q)\) y \(r=e^{-1/2}b_* /(1+b_*)\). Entonces
\[
p_{\mathrm I}(1)=2-\frac{2q}{3-2q},
\qquad p_{\mathrm{II}}(1)=2-\frac r{1-r}.
\]
Ambos programas producirían finalmente \(2\,\mathrm{mol}\) si se dejara reaccionar indefinidamente; ese límite no determina el producto al minuto.
